\chapter{後續學習方向}
\label{ch:where-next}

第 \ref{ch:introduction} 章結尾，你剛寫完人生第一支跑得動的程式，那時候書裡先打了個預防針：接下來會用
兩個方式反覆幫你建立 Python 的直覺，也會帶你去兩個地方，一個是讓你自己動手解決生活
裡的瑣事，一個是讓你看懂軟體這個世界怎麼運作。二十六章之後，這一章要做的事情，是回頭
兌現那個預告。

這一章分幾段推進：

\begin{itemize}
    \item 先回頭看，把你從一支只會印出一行字的程式，走到現在能處理資料檔案、串接真正的
    網路服務、交出一個有好幾個檔案的專案，這一路上兩個心智模型到底怎麼一步步落地，
    一件一件指給你看。
    \item 接著把兩個目標拆開來看，具體點出你在每個目標底下各自練成了哪些看得見、用得上
    的技能。
    \item 再把視角拉高一層，談真正能帶著走的東西，是幾個學習的習慣，讀錯誤訊息、查文件、
    遇到不會的東西先動手試一次，這些習慣以後遇到任何新語言、新工具都用得上。
    \item 然後把地圖攤開，讓你知道這本書刻意沒有走進去的幾片區域長什麼樣子，更深的
    Python 語言本身、更深的資料工具，還有幾種你以後可以自己找到的學習資源。
    \item 最後，依照你自己的背景，給一個具體一點的下一步建議。
\end{itemize}

這一章不會有新的片語、新的範例要你動手照著打，它的工作是幫你把已經學會的東西看清楚，
然後好好地把書闔上。

\section{回頭看：這趟旅程，你其實走了多遠}
\label{sec:journey-recap}

翻到這本書的最後一章之前，你大概不會沒事回頭翻前面二十六章，但現在正是這麼做的時候，
目的很單純：讓你老實對自己說一句，你真的走了這麼遠。這趟旅程從一行只會印出「哈囉，
世界！」的程式開始（\ref{ch:introduction} 章）；沒多久，你寫的程式第一次能自己做決定，依照某個條件走
哪一條分岔路（\ref{ch:conditionals} 章條件判斷）；再過不久，它學會了重複做一件事，不用你自己一行一行
複製貼上（\ref{ch:loops} 章迴圈）。中段，你的程式第一次打開一份真正的資料檔案，把裡面雜亂的內容
整理成看得懂的樣子，處理的對象從你自己手打的幾個數字，變成別人留下來、格式不一定
乾淨的原始資料（\ref{ch:csv} 章 CSV）。

再往後，你的程式碼本身開始變了樣子：四行迴圈加 append 縮成一行
comprehension（\ref{ch:comprehensions} 章）、一個函式學會「要一個才給一個」而不是先把
全部算完（\ref{ch:iterators-generators} 章）、一個函式敢在自己的主體裡呼叫自己
（\ref{ch:recursion} 章）。這些不是新功能，是同一件事寫得更像 Python 的樣子。

最後，它甚至跨出了你的電腦，主動去問另一個系統要資料，
再把回應接回自己的程式裡（\ref{ch:api} 章 API 串接）；而在 \ref{ch:capstone} 章，你把這一路
學到的東西全部放進同一個資料夾，交出了一個有好幾個檔案、有 \texttt{README}、有檢查腳本的
專案——那已經不是「一支程式」，是一個小小的軟體。這幾個里程碑，你是自己一行一行敲出來、
看著螢幕真的照你的話動起來的，這一段記憶，比任何一條語法規則都更值得留下來。

把這些具體的記憶收好之後，值得回頭看看，這一路你反覆撞見的到底是什麼。這本書一開始
就打過預防針：接下來會用兩個方式反覆幫你建立 Python 的直覺，現在是把這兩條線攤開
來看的時候。第一個是「名字牌模型」：\ref{ch:variables} 章第一次完整講給你聽，指定的意思是把一張名字
牌貼到一個物件上，跟把值複製進一個盒子完全是兩回事；之後你在 \ref{ch:lists} 章看 list 被兩個名字
同時指著、改一個兩個都動的時候撞見它一次，在 \ref{ch:functions} 章看函式怎麼「借」物件給參數的時候又
撞見一次，在 \ref{ch:tuples} 章看 tuple 為什麼不能被原地修改、名字牌模型跟盒子模型在這裡難得一致
的時候再一次，最後在 \ref{ch:advanced-functions} 章，你看見一個函式被回傳出來之後，居然還
記得出生時外層那張名字牌。

第二個是 EAFP：\ref{ch:files-exceptions} 章第一次完整鋪陳，遇到可能失敗的操作，
Python 比較習慣先做、錯了再處理，跟先檢查再動手的直覺剛好相反；你在 \ref{ch:dict-set} 章查字典之前
先埋過伏筆，在 \ref{ch:debugging} 章讀錯誤訊息、真正用 \texttt{try}/\texttt{except} 接住例外的時候，
親手用過它。

這兩條線在 \ref{ch:pythonic} 章已經合流過一次了——那一章把「以為 \texttt{b = a} 會複製」
跟「以為動手前應該先檢查」擺在一起，指出它們其實是同一個直覺，而且都猜錯了同一件事。
所以這裡不再重講一遍。 這裡只補一句那時候還不能說的話：那一章結束時，這兩個模型從
「用來解釋怪事的圖」變成了「你打字時會用到的規則」，而現在，隔了好幾章之後，你大概
已經不太會意識到自己在用它們了——這正是一個心智模型該有的下場。

如果你想幫這段回顧留下一個具體的動作，不妨真的翻回 \ref{ch:variables} 章你寫的第一個關於變數的小範例，
想想現在的你會怎麼重新寫一次它，會不會多包一個函式，讓它更整齊一點？或者翻回 \ref{ch:files-exceptions} 章
第一次寫的 \texttt{try}/\texttt{except}，想想現在的你，是不是已經很自然地先想到用它，
而不會先寫一堆 \texttt{if} 檢查？答案不用寫下來，你自己心裡有數就夠了，這正是這一路
上真正留下來的東西。

\section{兩個目標，你各自帶走了什麼技能}
\label{sec:goals-recap}

這本書一開始講過兩個並重的目標，這裡把它們拆開來看，具體點出你在每個目標底下各自練
成了哪些技能，一件一件指給你看，不停在「你進步了」這種空泛的講法上。

第一個目標，讓你能自己動手解決生活裡的瑣事。回頭數數看，你其實已經湊齊了一整套處理
資料的基本功：

\begin{itemize}
    \item \ref{ch:strings} 章教你用字串方法把一段雜亂的文字清乾淨。
    \item \ref{ch:lists} 章跟 \ref{ch:dict-set} 章教你用 list 跟 dict 把一批資料組織起來，\ref{ch:functions} 章教你把重複用得到的
    邏輯包成函式，不用每次都重寫一遍。
    \item \ref{ch:files-exceptions} 章跟 \ref{ch:csv} 章把這些串起來，讓你真的打得開一份檔案、讀得懂一份 CSV，處理完
    再存回去。
    \item \ref{ch:data-tools} 章讓你認識 numpy、pandas、matplotlib 這幾個常見的資料工具長什麼樣子，
    就算現在還不熟練，至少以後看到別人用這幾個名字，你知道它們在做什麼。
    \item \ref{ch:debugging} 章給你一套讀錯誤訊息、抓蟲的方法，讓「程式不動」這件事從一場災難變成
    一個可以按部就班解決的小問題。
    \item \ref{ch:regex} 章的正規表達式，讓你多一種在一大段文字裡揪出符合特定樣式內容的能力。
    \item \ref{ch:advanced-functions} 章到 \ref{ch:pythonic} 章這五章，讓同樣一件事開始寫得
    像 Python：把函式當成資料傳來傳去、用一行 comprehension 取代四行迴圈、讓資料「要一個
    才給一個」、遇到層層包住的東西時用遞迴拆開，以及一整套「別人一看就懂」的寫法慣例。
    這幾章沒有讓你多會做一件新的事，但讓你寫出來的東西，別人讀得下去。
    \item \ref{ch:automating} 章把這些技巧全部兜在一起，做出一個真正能跑的小專案，讓你體驗過一次
    「從一個模糊的需求，走到一支真的能用的程式」完整是什麼感覺，不再只是零散的片段
    練習。
\end{itemize}

這一路走下來，下次生活裡出現一件瑣碎、重複、要處理一堆資料的事情，你的第一反應會是
「這個我可以寫程式做」，而且真的知道從哪裡下手。

第二個目標，讓你看懂「軟體」這個世界怎麼運作。

\begin{itemize}
    \item \ref{ch:objects} 章跟 \ref{ch:inheritance} 章讓你知道物件、類別、繼承這些詞在講什麼，這是理解大部分現成
    程式庫、框架內部長相的基礎，就算你以後不自己寫大型系統，看得懂這套語彙也有幫助。
    \item \ref{ch:industry} 章一口氣把版本控制、前端、後端、資料庫、staging、production、backlog、
    sprint 這些工程師團隊平常掛在嘴邊的詞攤開講清楚，第一次讓你聽懂了他們平常在講的
    那些詞。
    \item \ref{ch:api} 章讓你親手讓自己的程式跟另一個系統打了一次交道，體會到「軟體跟軟體對話」
    實際上長什麼樣子，這件事光用聽的很難有感覺，只有自己打過一次 API 才會真的懂。
    \item \ref{ch:capstone} 章讓你把 \ref{ch:industry} 章那些只是「聽懂」的詞，親手做了一次：
    多個檔案怎麼分、\texttt{requirements.txt} 寫什麼、\texttt{README} 該交代什麼、
    改壞了要怎麼發現。 那一章交出來的東西，形狀已經跟工程師交出來的東西是同一種了。
\end{itemize}

這兩個目標從頭到尾沒有誰讓位給誰，你這一路走來，兩邊都真的碰過一次，兩個
都是走到今天、真正握在手上的東西。

\section{把方法帶走：比語法更重要的，是幾個習慣}
\label{sec:habits-to-keep}

技巧會忘、語法會查，但有幾件事，你已經在不知不覺中練成了習慣，而且以後不管換到哪
一種語言、哪一套工具，都還是用得上。

\textbf{讀錯誤訊息的習慣。} \ref{ch:introduction} 章你第一次看到一片紅字的時候，這本書就跟你說過：由下往
上讀，先看最後一行的錯誤類型跟訊息，那才是真正的線索；\ref{ch:debugging} 章你又把這套習慣練得更熟。
這件事情不限於 Python，以後不管你換到哪一種語言、用哪一套工具，錯誤訊息的樣子會不
一樣，但「看到紅字先別慌，由下往上找線索」這套習慣，你已經帶在身上了，搬到哪裡都用
得上。

\textbf{查文件的習慣。} \ref{ch:modules} 章這本書刻意沒有要你把每個標準函式庫的函式都背下來，改成
帶你練習「查文件」這件事本身，因為沒有人真的把整個 Python 標準函式庫背在腦子裡，連
寫了很多年程式的人也不例外。以後遇到書裡沒教過的模組、沒教過的套件，官方文件或套件
自己附的說明，永遠是你第一個該去的地方，不用繼續憑猜。

\textbf{遇到不會的東西先動手試一次的習慣。} \ref{ch:introduction} 章教你用互動模式「隨手驗算」的時候，這
本書就已經在示範這個習慣了，不確定某段語法會怎麼動，與其在腦子裡空想，不如直接打開
來試一次，答案馬上就出來了。

這件事其實也是 EAFP 精神的另一種面貌，EAFP 教你的是
「先做、錯了再處理」的程式寫法，但拉高一層來看，它其實也是一種面對陌生事物的態度：
遇到沒學過的東西，先動手碰一次，錯了再回頭修正，往往比坐在那裡把所有可能情況都想過
一遍更快找到答案。這裡是把 EAFP 這個詞往外延伸成一種學習態度，用同一種精神套用在
「學新東西」這件更大的事情上，一樣說得通，它原本在程式碼裡的意思還是那個意思。

\textbf{把大問題拆成小步驟的習慣。} \ref{ch:automating} 章那個整合專案，如果從頭到尾當成一件事來想，
很容易讓人不知道從何下手；但只要把它拆成幾個小任務，一個一個各個擊破，最後接起來就
是完整的成品。

這件事其實從第一次寫函式（\ref{ch:functions} 章）就開始練了，函式本身就是「把一段完整
邏輯包成一個獨立、可以單獨測試的小塊」的做法。以後遇到任何看起來很龐大的專案或任務，
先問自己「這件事可以拆成哪幾個更小的步驟」，往往比想著「怎麼一次寫出整個解法」更容易
真正開始動手。

\textbf{以後如果要用 AI 工具輔助寫程式或除錯。} 這本書沒有教你怎麼用 ChatGPT 之類的
工具生成程式碼，但值得老實說一句：這幾個你已經練出來的習慣，讀錯誤訊息、查文件、
動手先試、把問題拆小，以後一樣用得上，甚至是你判斷「AI 給的答案到底對不對」的基礎。

AI 工具能幫你很快生出一段程式碼，但那段程式碼有沒有真的做了你要的事，有沒有藏著一個
你沒發現的錯誤，還是得靠你自己讀得懂、查得到、試得出來。這本書真正想讓你帶走的，是
有能力判斷 AI 給的東西靠不靠得住，不是把以後所有的事都交給它。

\section{地圖上還有更大的世界}
\label{sec:beyond-this-book}

前面幾節整理的是你已經走過的路；這一節要做的事情不一樣，把地圖攤開，指給你看這本書
刻意沒有走進去的幾片區域長什麼樣子。目的不是要你現在就衝進去，只是讓你知道它們在那裡，
以後想去的時候認得路。

\textbf{更深的 Python 語言本身。} 這本書教的是 Python 能做到的一部分，全貌還大得多。
有幾件事，你在這本書裡碰到的是入門的那一層，往下還有好幾層；也有幾個詞，你完全沒碰過，
但以後跟程式圈的人聊起 Python，大概率會聽到：

\begin{itemize}
    \item \texttt{async}／\texttt{await}，讓程式在等待網路回應這種「發呆時間」裡，先去
    做別的事，不用乾等。第 \ref{ch:api} 章你打 API 的時候，如果一次要打很多個網址，
    就會用得上這個。
    \item \textbf{自己寫一個 context manager}。第 \ref{ch:pythonic} 章你已經知道
    \texttt{with} 在解決什麼問題了，但那一章只教你怎麼用別人做好的。等你學會怎麼替自己的
    東西加上「進去做什麼、出來做什麼」這一層，很多重複的收尾動作就可以收起來。
    \item \textbf{decorator 更完整的樣子}。第 \ref{ch:advanced-functions} 章教的是最
    基本的形狀，真正的專案裡你會看到帶參數的 decorator、疊好幾層的 decorator，以及
    \texttt{functools} 裡幫你把這些寫對的工具。
    \item \textbf{型別檢查工具}。第 \ref{ch:advanced-functions} 章說過，型別標註不會
    影響程式執行，Python 不會幫你檢查——但有另外一類工具會。在多人維護的專案裡，它們
    能在程式跑起來之前就抓出「這裡傳錯東西了」。
    \item \textbf{測試框架}。第 \ref{ch:capstone} 章你用幾個 \texttt{assert} 寫過一支
    檢查腳本，那是測試最素樸的樣子。真正的專案會用專門的測試工具，可以一次跑幾百個檢查、
    告訴你哪一個掛了、還能算出你的程式碼有多少比例被測到。
    \item \textbf{把程式碼分裝成套件}。第 \ref{ch:modules} 章教的自訂模組，前提是檔案跟
    呼叫它的程式放在同一個資料夾。專案再大一點就會需要子資料夾、需要讓別人可以
    \texttt{pip install} 你寫的東西，那是另一套規則。
\end{itemize}

這幾個詞，這裡只給你一句話的印象，不寫成可以動手照打的片語
或範例，這是這一節跟這本書其他每一節最大的不同。

這些詞具體會落在哪裡，可以看一個實際
的例子：像《Fluent Python》這樣的書，會把 iterator 與 generator、context managers、
\texttt{async}／\texttt{await} 這些放進同一個部分，書裡稱為「Control Flow」，用完整
一個部分的份量系統性地教——你在第 \ref{ch:iterators-generators} 章學到的 \texttt{yield}，
在那本書裡只是那一整個部分的開頭。那本書的作者甚至在一開頭就老實提醒讀者：「如果你才剛開始
學 Python，這本書會很難跟上，甚至可能讓你誤以為每一支 Python 程式都得用上特殊方法跟
metaprogramming 技巧。」這句話送給正在看這本書最後一章的你剛好合適：那本書現在還不
適合你，但它在地圖上，以後你想往「更深的 Python 語言本身」走的時候，認得路就好。

\textbf{更嚴謹的開發習慣。} 這本書裡的每支程式，寫完能跑、輸出結果對，任務就算完成。
但一支程式如果要讓很多人一起維護、要長期跑在正式環境裡，光是「能跑」還不夠。
\texttt{testing}（測試），是指另外寫一段程式碼，專門用來自動檢查你原本那段程式碼
在各種情況下是不是都給出正確的結果，這樣以後改動程式碼時，只要重新跑一次測試，就能
立刻知道有沒有不小心弄壞了原本能動的部分，不用每次都靠人工從頭試一遍。

前面提到的
\texttt{type hints}，也常常被歸在這一類「讓程式碼更可靠」的做法裡。這兩件事這本書
完全沒有觸碰，因為它們的價值要在程式碼變大、變得需要跟別人協作之後才真正顯現出來，
現在的你還不需要，但以後往這個方向走的時候，知道有這兩個詞可以查，會省下不少繞路
的時間。

\textbf{更大的應用範圍。} \ref{ch:api} 章你打 API 的時候，是你的程式主動去問別人的伺服器要
資料，扮演「使用者」的角色。

如果反過來，你想讓自己寫的程式變成別人可以在瀏覽器裡
打開、或是別的程式可以打過來要資料的服務，你會用到「網頁框架」這一類工具，它們幫你
把接收請求、整理回應這些底層雜事先打包好，讓你可以專心寫「收到請求之後要做什麼」這
段邏輯。這一塊跟 \ref{ch:industry} 章講的前端、後端、API 這些詞是同一條線上的東西，只是這本書只帶
你從「使用者」的角度碰過一次，還沒帶你走到「自己就是提供服務的那一方」這一步。

\textbf{更深的資料工具。} \ref{ch:data-tools} 章你已經跟 pandas 打過照面，用它讀 CSV、選欄、篩出符合
條件的資料列。但那裡刻意留了兩個詞沒有教：

\begin{itemize}
    \item \texttt{groupby} 做的事情是「把資料依某個欄位分組，再對每組分別算統計量」，
    例如「這份氣溫資料裡，每個城市個別的平均溫度是多少」；\ref{ch:data-tools} 章你如果想做這件事，得
    自己寫迴圈或條件判斷一步步湊出來，\texttt{groupby} 能把整段邏輯壓縮成一行。
    \item \texttt{merge} 做的是「把兩份資料依共同的欄位接在一起」，例如你手上有一份
    訂單資料、一份客戶資料，兩邊都有「客戶編號」這個欄位，\texttt{merge} 能幫你把兩份
    資料合成一份，不用自己一筆一筆對照著找。
\end{itemize}

這兩個詞，這裡一樣只給你一句話的印象，不會像前面章節那樣寫成完整的片語跟
範例。

\section{持續學習的資源}
\label{sec:learning-resources}

如果你想繼續往下走，不管是想多學一點寫給自己用的小工具，還是真的想認真往這個方向
發展，你不用一個人埋頭摸索。市面上跟網路上能找到的資源大致可以分成幾類，各有各的
用途，挑適合自己現在階段的就好，不用每一種都追。

第一類是書。

\begin{itemize}
    \item 有的書走「系統性參考」路線，把一個主題的來龍去脈講得很完整，適合你已經
    對某個方向有點興趣、想深入紮根的時候翻。
    \item 有的書走「動手做專案」路線，帶你一步步做出幾個完整的小成品，適合你想累積
    實作經驗、把書上的東西真正變成自己能重現的能力的時候用。
    \item 還有一種比較像速查手冊，平常不會從頭到尾讀，但寫程式寫到一半忘記某個語法時，
    翻兩下就能找到答案，適合放在手邊當工具書用。
\end{itemize}

第二類是線上課程。

\begin{itemize}
    \item 有的是一支支影片，老師帶著你邊講邊打字，適合喜歡「有人在旁邊帶」的學習方式。
    \item 有的是互動式的練習平台，網頁上直接寫程式碼、馬上看到對不對，適合想邊學邊
    動手、不想只是看的人。
    \item 還有一種比較像結構完整的大學課程搬到線上，會有作業、有進度安排，適合想要
    有人幫你排好學習路徑、按表操課的人。
\end{itemize}

這幾種形式教的內容常常重疊，選哪
一種主要看你自己比較適應哪種學習節奏，不是哪一種天生比較高級。

第三類是練習網站。

\begin{itemize}
    \item 有的是一題一題的程式解題平台，出一個小問題、你寫程式解掉它，適合拿來磨
    基本功、練習「看到問題馬上想到怎麼用程式解」的反射。
    \item 有的偏向真實世界的小專案練習，給你一個比較貼近實際工作場景的任務，讓你
    練習把好幾個技巧兜在一起用。
    \item 還有一種是參與別人已經寫好、公開讓大家一起維護的專案，從讀懂別人的程式碼
    開始，慢慢學會怎麼跟其他人一起改動同一份程式碼，這也是很多工程師真正精進技巧的
    地方。
\end{itemize}

這一節刻意沒有點名任何一本書或任何一個網站，因為書會改版、網站會關站，但上面這幾類
資源的分工不太會變。 如果你現在就想要具體的書名，書末的〈參考資料〉列了這本書寫作時
實際參考的七本 Python 入門書，每一本都註明了適合在什麼時候讀，以及這本書的哪幾章借用
過它。

除了這些，Python 官方文件本身就寫得很完整，大部分你以後會用到的模組、套件，第一手
的說明都在那裡，\ref{ch:modules} 章已經帶你練過怎麼讀這種文件。

如果你住的城市附近剛好有 Python
的使用者社群，他們通常對新手很友善，不用怕自己「還不夠格」加入；找不到實體社群的話，
線上的討論區、留言板一樣有很多人在，遇到問題丟出去問，通常都會有人願意回你。這幾類
資源不用照順序一個一個走完，遇到卡關的地方，哪一種能幫上忙就先用哪一種。

\section{接下來，你是誰，就從哪裡先走}
\label{sec:next-steps-by-reader}

前面幾節給的建議比較通用，這一節具體一點，依照你自己的背景，給一個下一步的方向。你
不用完全對號入座，這幾種只是常見的起點，挑一個離你目前狀況最近的參考就好。

\begin{itemize}
    \item \textbf{國中生或高中生：} 現在最值得做的事情，不是急著衝去學更深的東西，是
    把這本書已經學過的東西多用幾次，變成真正內化的手感。試著幫自己寫一個小工具，整理
    作業進度、算一下成績、幫喜歡的遊戲寫個小腳本，題目越貼近你自己的生活越好。如果
    學校或住的城市附近有程式相關的社團或營隊，去看看也無妨，重點是持續動手，比急著
    往前衝更重要。
    \item \textbf{文組背景的學生，或原本從來沒接觸過理工訓練的人：} 這本書想讓你確認
    的一件事是：你不需要轉去讀資訊工程，也能真的把 Python 用起來。可以把它當成原本
    專業裡多出來的一項工具，用來整理研究資料、自動化一些重複的文書工作，或是單純
    因為好奇多一項技能。下一步不用是「換一個更難的領域」，可以是「把 Python 用回你
    原本在乎的那件事情上」。
    \item \textbf{已經在職場工作，尤其是 PM、營運、行銷這類常常要跟工程師對接、但
    自己不寫程式的角色：} \ref{ch:industry} 章跟 \ref{ch:api} 章大概是這本書對你最直接有用的兩章。下一步可以
    試著在下次開會前，先把想問工程師的問題用這本書教過的詞彙組織一遍，你會發現溝通
    明顯順暢很多。如果手邊剛好有重複性高的報表整理或資料彙整工作，也很值得試著自己
    寫一支小程式接手，不需要變成工程師，只要能看懂、能跟工程師討論、偶爾自己動手
    解決一點小麻煩，這本書設定的目標就已經達成了。
\end{itemize}

不管你屬於哪一種背景，這幾條路徑最後其實會匯合到同一個地方：持續用、持續碰、持續在
真實情境裡解決一件具體的小問題，比一口氣把地圖上剩下的每個區域都學完更重要。

這本書開頭第 \ref{ch:introduction} 章，你剛寫完人生第一支跑得動的程式，那是第一章的尾聲，而且已經有一支
自己跑得動的程式。

現在，你在整本書的尾聲，你手上有的，不只是一些程式碼片段，是兩種
看世界的方式：一種讓你看到生活裡的瑣事，會反射性想「這個我可以寫程式做」；另一種讓
你看懂「軟體」這個世界平常怎麼運作，不再被術語唬住。這兩件事，從第一頁開始就沒有誰
比較重要，現在，它們都真的是你的了。
